\documentclass[11pt,a4paper]{article}
\usepackage[T1]{fontenc}
\usepackage[utf8]{inputenc}
\usepackage{amsmath,amssymb,xcolor}
\usepackage[margin=25mm]{geometry}
\usepackage[colorlinks=true,urlcolor=blue]{hyperref}
\definecolor{nexorablue}{HTML}{2459B5}
\setlength{\parindent}{0pt}
\setlength{\parskip}{7pt}
\setlength{\emergencystretch}{2em}
\renewcommand{\familydefault}{\sfdefault}
\begin{document}

{\small\bfseries\color{nexorablue} NEXORA CPGE}\par
{\small MPSI / NEX-MES-06 / Corrigé}\par
{\LARGE\bfseries Un zéro partagé par deux mesures}\par
\vspace{0.25em}\hrule\vspace{0.6em}

\section*{Modèle à erreur commune}
Le décalage s’annule dans la différence, donc $u(d)=\sqrt{0{,}02^2+0{,}03^2}\simeq0{,}036\,\mathrm{cm}$.

\section*{Comparaison}
En supposant à tort l’indépendance des mesures totales, on ajouterait deux fois $0{,}05^2$ à la variance de la différence. Cela donnerait environ $0{,}079\,\mathrm{cm}$. Le capteur est le même dans les deux lectures, d’où la corrélation positive.


\end{document}
